\documentclass[aspectratio=169,10pt]{beamer}
\makeatletter\def\input@path{{../../estilo/}{../estilo/}{estilo/}}\makeatother
\usepackage{devmobile}
\modulo[Kotlin]{06}{Kotlin Assíncrono: Coroutines e Flow}{Thread principal, funções suspensas, dispatchers, concorrência estruturada, Flow e StateFlow}

\begin{document}
\capa

\begin{frame}{Objetivos do módulo}
  \begin{conceito}[Ao final deste módulo você será capaz de]
    \begin{itemize}
      \item Explicar por que operações lentas não podem rodar na \emph{thread} principal.
      \item Escrever funções \cod{suspend} e iniciar coroutines com \cod{launch} e \cod{async}.
      \item Escolher o \emph{dispatcher} adequado com \cod{withContext}.
      \item Aplicar concorrência estruturada, cancelamento e tratamento de exceções.
      \item Produzir e consumir fluxos de dados com \cod{Flow}.
      \item Expor estado observável com \cod{StateFlow} e eventos com \cod{SharedFlow}.
    \end{itemize}
  \end{conceito}
\end{frame}

\sumario

% =============================================================================
%  Fundamentos compartilhados --- Módulo 06: concorrência e assincronia
%  Incluído pelas trilhas Kotlin e Flutter com % =============================================================================
%  Fundamentos compartilhados --- Módulo 06: concorrência e assincronia
%  Incluído pelas trilhas Kotlin e Flutter com % =============================================================================
%  Fundamentos compartilhados --- Módulo 06: concorrência e assincronia
%  Incluído pelas trilhas Kotlin e Flutter com \input{../fundamentos/f06-concorrencia}
% =============================================================================
\section{Fundamentos: concorrência e assincronia}

\begin{frame}{Programas, processos e \emph{threads}}
  \begin{columns}[c]
    \begin{column}{0.47\textwidth}
      \centering
      \begin{tikzpicture}[every node/.style={font=\tiny}]
        \draw[azulpetroleo,line width=0.8pt,rounded corners=4pt,fill=cinzafundo] (0,0) rectangle (5.6,3.2);
        \node[rotulo,anchor=north west] at (0.1,3.15) {\textbf{processo do app}};
        \node[caixaambar,minimum width=5.1cm,minimum height=5mm,font=\tiny] at (2.8,2.45) {memória compartilhada (\emph{heap}): objetos do app};
        \foreach \x/\n in {0.95/{principal (UI)},2.8/{rede},4.65/{cálculo}}{
          \node[caixaescura,minimum width=1.6cm,minimum height=11mm,font=\tiny,align=center] at (\x,1.15) {\emph{thread}\\ \n\\[1pt] \textcolor{ambarclaro}{pilha própria}};
        }
        \node[caixa,minimum width=5.6cm,minimum height=5mm,font=\tiny] (so) at (2.8,-0.55) {sistema operacional: escalonador};
        \foreach \x in {0.7,2.1,3.5,4.9}{\node[caixa,minimum width=1.1cm,minimum height=4mm,font=\tiny] at (\x,-1.3) {núcleo};}
      \end{tikzpicture}
    \end{column}
    \begin{column}{0.51\textwidth}
      \begin{definicao}[Processo e \emph{thread}]
        \textbf{Processo}: programa em execução, com memória própria e isolada dos outros apps. \textbf{\emph{Thread}}: linha de execução dentro do processo; as \emph{threads} dividem a memória do processo, mas cada uma tem sua pilha.
      \end{definicao}
      \begin{pordentro}[O escalonador]
        O sistema operacional distribui as \emph{threads} pelos núcleos e alterna entre elas milhares de vezes por segundo (troca de contexto). Cada \emph{thread} custa memória e tempo de troca: não se cria uma por tarefa.
      \end{pordentro}
      \begin{analogia}
        O processo é a cozinha; as \emph{threads} são os cozinheiros. Eles dividem a despensa, e é aí que começam as brigas pelo mesmo ingrediente.
      \end{analogia}
    \end{column}
  \end{columns}
\end{frame}

\begin{frame}{Concorrência e paralelismo}
  \begin{columns}[c]
    \begin{column}{0.47\textwidth}
      \centering
      \begin{tikzpicture}[x=0.55cm,y=0.5cm,every node/.style={font=\tiny}]
        \node[rotulo,anchor=east] at (-0.2,2) {1 núcleo};
        \foreach \x/\cor in {0/petroleo,1.5/ambar,3/petroleo,4.5/ambar,6/petroleo,7.5/ambar}{\fill[\cor!80] (\x,1.6) rectangle ++(1.4,0.8);}
        \node[rotulo] at (4.5,3.1) {\textbf{concorrência}: as tarefas se alternam};
        \node[rotulo,anchor=east] at (-0.2,0) {núcleo 1};
        \fill[petroleo!80] (0,-0.4) rectangle ++(8.9,0.8);
        \node[rotulo,anchor=east] at (-0.2,-1.1) {núcleo 2};
        \fill[ambar!80] (0,-1.5) rectangle ++(8.9,0.8);
        \node[rotulo] at (4.5,-2.3) {\textbf{paralelismo}: as tarefas rodam juntas};
        \draw[->,cinzamedio] (0,-2.9) -- (9.3,-2.9) node[rotulo,right]{tempo};
        \fill[petroleo!80] (0.4,4) rectangle ++(0.5,0.4); \node[rotulo,anchor=west] at (1,4.2) {tarefa A};
        \fill[ambar!80] (3.4,4) rectangle ++(0.5,0.4); \node[rotulo,anchor=west] at (4,4.2) {tarefa B};
      \end{tikzpicture}
      \raggedright
      \begin{pordentro}[Chips de celular]
        Combinam núcleos de alto desempenho e núcleos econômicos (big.LITTLE): o sistema escolhe onde cada \emph{thread} roda para poupar bateria.
      \end{pordentro}
    \end{column}
    \begin{column}{0.51\textwidth}
      \begin{definicao}[Concorrência]
        Várias tarefas em andamento no mesmo período, que podem se alternar. É um jeito de \textbf{estruturar} o programa.
      \end{definicao}
      \begin{definicao}[Paralelismo]
        Várias tarefas executando \textbf{no mesmo instante}, em núcleos diferentes. É um jeito de \textbf{executar}.
      \end{definicao}
      \begin{historia}[Na frase de Rob Pike (2012)]
        Concorrência é lidar com muitas coisas ao mesmo tempo; paralelismo é fazer muitas coisas ao mesmo tempo.
      \end{historia}
    \end{column}
  \end{columns}
\end{frame}

\begin{frame}{Tarefas de I/O e tarefas de CPU}
  \scriptsize
  \begin{tabularx}{\textwidth}{@{}l L L l l@{}}
    \toprule
    \cabecalho{Tipo} & \cabecalho{O que limita} & \cabecalho{Exemplos} & \cabecalho{Kotlin} & \cabecalho{Dart} \\
    \midrule
    I/O & espera por rede, disco ou banco & requisição HTTP, consulta SQL, ler arquivo & \cod{suspend} + \cod{Dispatchers.IO} & \cod{Future} + \cod{await} \\
    CPU & cálculo no processador & decodificar imagem, JSON gigante, criptografia & \cod{Dispatchers.Default} & \cod{Isolate.run} \\
    Tempo & o relógio & atraso, repetição, tempo limite & \cod{delay}, \cod{withTimeout} & \cod{Future.delayed}, \cod{Timer} \\
    \bottomrule
  \end{tabularx}
  \normalsize
  \vspace{4pt}
  \begin{columns}[T]
    \begin{column}{0.49\textwidth}
      \begin{definicao}[Bloqueante \emph{versus} não bloqueante]
        Uma chamada \textbf{bloqueante} prende a \emph{thread} até terminar. Uma \textbf{não bloqueante} registra o pedido e libera a \emph{thread}; o resultado chega depois, por notificação do sistema.
      \end{definicao}
    \end{column}
    \begin{column}{0.49\textwidth}
      \begin{conceito}[A ideia-chave do módulo]
        \textbf{Esperar não precisa de \emph{thread}}: enquanto a resposta não chega, a mesma \emph{thread} atende outros eventos. Já o \textbf{cálculo pesado precisa de um núcleo} de verdade, fora da \emph{thread} da interface.
      \end{conceito}
    \end{column}
  \end{columns}
\end{frame}

\begin{frame}{O laço de eventos em cada plataforma}
  \begin{columns}[T]
    \begin{column}{0.55\textwidth}
      \begin{definicao}[Laço de eventos (\emph{event loop})]
        Uma \emph{thread} que repete sem parar: pega o próximo evento da fila, executa o código associado até o fim e volta à fila. Toda interface gráfica funciona assim.
      \end{definicao}
      \scriptsize
      \begin{tabularx}{\linewidth}{@{}l L@{}}
        \toprule
        \cabecalho{Plataforma} & \cabecalho{Mecanismo na \emph{thread} principal} \\
        \midrule
        Android & \cod{Looper} + \cod{MessageQueue} + \cod{Handler} \\
        iOS & \cod{RunLoop} e fila principal do GCD \\
        Dart & fila de microtarefas e fila de eventos \\
        Web & laço de eventos do navegador (tarefas e microtarefas) \\
        \bottomrule
      \end{tabularx}
    \end{column}
    \begin{column}{0.43\textwidth}
      \begin{pordentro}[Dart: duas filas]
        Antes de cada evento, o Dart esvazia a fila de \textbf{microtarefas} (continuações curtas de \cod{Future}s). Desde o Flutter 3.29, no Android e no iOS, o código Dart roda na própria \emph{thread} principal do app.
      \end{pordentro}
      \begin{atencao}[Um evento lento atrasa todos]
        Quadros se perdem (\emph{jank}) e, no Android, sem resposta a um toque por cerca de 5\,s aparece o diálogo ANR.
      \end{atencao}
    \end{column}
  \end{columns}
\end{frame}

\begin{frame}{Da pirâmide de \emph{callbacks} ao \cod{async}/\cod{await}}
  \centering
  \begin{tikzpicture}[every node/.style={font=\tiny}]
    \draw[seta] (0,0) -- (13.4,0);
    \foreach \x/\t/\d in {1.2/\emph{callbacks}/{funções chamadas\\ no fim da tarefa},4.4/\emph{Futures} e \emph{Promises}/{objeto que representa\\ o resultado futuro},7.8/{\cod{async}/\cod{await}}/{C\# (2012), Dart (2015),\\ JavaScript (2017)},11.4/concorrência estruturada/{Kotlin (2018):\\ escopos e cancelamento}}{
      \fill[ambar] (\x,0) circle (2pt);
      \node[font=\scriptsize\bfseries,text=azulpetroleo,above=1mm] at (\x,0) {\t};
      \node[rotulo,below=1mm] at (\x,0) {\d};
    }
  \end{tikzpicture}
  \raggedright
  \begin{columns}[T]
    \begin{column}{0.49\textwidth}
      \begin{definicao}[\emph{Future} (ou \emph{Promise})]
        Objeto que representa um resultado que ainda não existe. Começa \textbf{pendente} e termina \textbf{com valor} ou \textbf{com erro}. Quem o recebe pode esperar sem bloquear a \emph{thread}.
      \end{definicao}
      \centering
      \begin{tikzpicture}[x=0.5cm,y=0.45cm,every node/.style={font=\tiny}]
        \node[rotulo,anchor=east] at (-0.2,0.4) {\emph{thread}};
        \fill[petroleo!80] (0,0) rectangle ++(2,0.8); \node[white] at (1,0.4) {trecho 1};
        \fill[cinzaclaro] (2,0) rectangle ++(3,0.8); \node at (3.5,0.4) {outros eventos};
        \fill[petroleo!80] (5,0) rectangle ++(2,0.8); \node[white] at (6,0.4) {trecho 2};
        \fill[cinzaclaro] (7,0) rectangle ++(2.4,0.8); \node at (8.2,0.4) {outros};
        \fill[petroleo!80] (9.4,0) rectangle ++(1.6,0.8); \node[white] at (10.2,0.4) {fim};
        \foreach \x in {2,7}{\draw[ambar,line width=1pt] (\x,-0.2) -- (\x,1.1); \node[rotulo,above] at (\x,1.1) {\cod{await}};}
      \end{tikzpicture}
    \end{column}
    \begin{column}{0.49\textwidth}
      \begin{pordentro}[O que acontece em um \cod{await}]
        A função pausa em um \textbf{ponto de suspensão} e devolve a \emph{thread}; quando o resultado chega, continua de onde parou, com as variáveis locais preservadas. O compilador do Kotlin transforma cada função \cod{suspend} em uma máquina de estados que recebe uma \cod{Continuation}.
      \end{pordentro}
      \begin{dica}
        O ganho é escrever código assíncrono com cara de sequencial: \cod{try}/\cod{catch} e laços funcionam como sempre.
      \end{dica}
    \end{column}
  \end{columns}
\end{frame}

\begin{frame}{Concorrência estruturada}
  \begin{columns}[c]
    \begin{column}{0.45\textwidth}
      \centering
      \begin{tikzpicture}[every node/.style={font=\tiny}]
        \node[caixaescura,minimum width=3.4cm,font=\tiny] (e) at (0,0) {escopo da tela};
        \node[caixa,minimum width=1.6cm,font=\tiny] (a) at (-0.95,-1) {carregar lista};
        \node[caixa,minimum width=1.6cm,font=\tiny] (b) at (0.95,-1) {ouvir sensor};
        \node[caixaambar,minimum width=1.6cm,font=\tiny] (c) at (0.95,-2) {ler histórico};
        \draw[seta] (e) -- (a); \draw[seta] (e) -- (b); \draw[seta] (b) -- (c);
        \draw[pontilhado] (-2,0.5) rectangle (2,-2.5);
        \node[rotulo,text width=4.2cm] at (0,-3) {a tela fecha $\Rightarrow$ o escopo termina $\Rightarrow$ as tarefas filhas são canceladas};
      \end{tikzpicture}
      \begin{historia}
        Assim como a programação estruturada aposentou o \cod{goto}, este modelo aposenta tarefas soltas. A ideia ganhou força em 2018, com um ensaio de Nathaniel J. Smith, e o Kotlin a adotou no mesmo ano.
      \end{historia}
    \end{column}
    \begin{column}{0.53\textwidth}
      \begin{definicao}[Concorrência estruturada]
        Toda tarefa concorrente nasce dentro de um \textbf{escopo} e não vive mais que ele: quando o escopo termina, as tarefas filhas são canceladas.
      \end{definicao}
      \begin{conceito}[Garantias]
        Nada fica rodando depois que a tela fecha; o cancelamento chega aos filhos; erros dos filhos chegam ao pai.
      \end{conceito}
      \begin{atencao}[E no Dart?]
        \cod{Future} e \cod{Stream} não seguem esse modelo: cancele inscrições e temporizadores no \cod{dispose} e confira \cod{mounted} após um \cod{await}.
      \end{atencao}
    \end{column}
  \end{columns}
\end{frame}

\begin{frame}{Condições de corrida e impasses}
  \begin{columns}[T]
    \begin{column}{0.47\textwidth}
      \centering
      \begin{tikzpicture}[x=0.7cm,y=0.6cm,every node/.style={font=\tiny}]
        \tikzset{passo/.style={caixa,minimum height=4.5mm,minimum width=1.2cm,inner sep=2pt,font=\tiny}}
        \node[rotulo,anchor=east] at (-0.1,1) {\emph{thread} A};
        \node[rotulo,anchor=east] at (-0.1,0) {\emph{thread} B};
        \node[passo] at (0.9,1) {lê 5};
        \node[passo] at (2.6,0) {lê 5};
        \node[passo,caixaambar,font=\tiny] at (4.3,1) {grava 6};
        \node[passo,caixaambar,font=\tiny] at (6,0) {grava 6};
        \draw[->,cinzamedio] (0,-0.7) -- (7,-0.7) node[rotulo,right]{t};
        \node[rotulo,text=vermelhoerro,text width=4.8cm] at (3.4,-1.3) {duas somas de 1 a partir de 5 deveriam dar 7: uma atualização se perdeu};
      \end{tikzpicture}
      \begin{definicao}[Condição de corrida]
        Erro em que o resultado depende da ordem imprevisível em que tarefas concorrentes acessam um dado compartilhado.
      \end{definicao}
      \begin{atencao}[Impasse (\emph{deadlock})]
        Duas tarefas, cada uma segurando um recurso e esperando o da outra: ninguém avança. Adquira travas sempre na mesma ordem.
      \end{atencao}
    \end{column}
    \begin{column}{0.51\textwidth}
      \scriptsize
      \begin{tabularx}{\linewidth}{@{}l L@{}}
        \toprule
        \cabecalho{Proteção} & \cabecalho{Como funciona} \\
        \midrule
        Exclusão mútua & uma tarefa por vez na região crítica (\cod{Mutex}) \\
        Operação atômica & ler e gravar em um passo indivisível (\cod{AtomicInteger}, \cod{StateFlow.update}) \\
        Confinamento & só uma \emph{thread} toca o dado (a principal, por exemplo) \\
        Imutabilidade & ninguém altera: cria-se uma cópia nova \\
        Mensagens & cada tarefa tem seus dados e troca cópias \\
        \bottomrule
      \end{tabularx}
      \normalsize
      \begin{pordentro}[Corrida lógica com uma \emph{thread} só]
        Entre dois \cod{await}, outro evento pode mudar o estado. Duas buscas disparadas em sequência podem responder fora de ordem, e a mais antiga sobrescrever a mais nova.
      \end{pordentro}
    \end{column}
  \end{columns}
\end{frame}

\begin{frame}{Memória compartilhada ou troca de mensagens}
  \begin{columns}[c]
    \begin{column}{0.43\textwidth}
      \centering
      \begin{tikzpicture}[every node/.style={font=\tiny}]
        \node[caixa,minimum width=1.3cm,font=\tiny] (a) at (0,0) {\emph{thread} A};
        \node[caixa,minimum width=1.3cm,font=\tiny] (b) at (2.6,0) {\emph{thread} B};
        \node[caixaambar,minimum width=1.6cm,font=\tiny] (o) at (1.3,-1.2) {objeto + trava};
        \draw[seta] (a) -- (o); \draw[seta] (b) -- (o);
        \node[rotulo] at (1.3,0.6) {\textbf{memória compartilhada}};
        \node[caixaescura,minimum width=1.3cm,minimum height=9mm,font=\tiny] (c) at (0,-3) {isolate A\\ memória A};
        \node[caixaescura,minimum width=1.3cm,minimum height=9mm,font=\tiny] (d) at (2.6,-3) {isolate B\\ memória B};
        \draw[setaambar] ([yshift=1.5mm]c.east) -- node[rotulo,above]{mensagem} ([yshift=1.5mm]d.west);
        \draw[setaambar] ([yshift=-1.5mm]d.west) -- node[rotulo,below]{resposta} ([yshift=-1.5mm]c.east);
        \node[rotulo] at (1.3,-2.1) {\textbf{troca de mensagens}};
      \end{tikzpicture}
    \end{column}
    \begin{column}{0.55\textwidth}
      \scriptsize
      \begin{tabularx}{\linewidth}{@{}l L L@{}}
        \toprule
         & \cabecalho{Memória compartilhada} & \cabecalho{Troca de mensagens} \\
        \midrule
        Exemplos & \emph{threads} da JVM (Kotlin, Java) & \emph{isolates} do Dart, atores, Web Workers \\
        Comunicação & ler e gravar o mesmo objeto & enviar cópias de dados \\
        Risco & corridas e impasses & custo de copiar dados grandes \\
        Proteção & travas, atômicos, confinamento & o próprio modelo isola \\
        \bottomrule
      \end{tabularx}
      \normalsize
      \begin{analogia}
        Um quadro branco compartilhado exige regras para ninguém apagar o do colega. Colegas que trocam bilhetes nunca brigam pelo quadro, mas gastam tempo escrevendo.
      \end{analogia}
    \end{column}
  \end{columns}
\end{frame}

\begin{frame}{Fluxos de dados: frios, quentes e contrapressão}
  \begin{columns}[T]
    \begin{column}{0.55\textwidth}
      \begin{definicao}[Fluxo (\emph{stream})]
        Sequência de valores que chegam ao longo do tempo: leituras de sensor, mensagens, mudanças no banco. Segue o padrão \textbf{Observador}: quem produz notifica quem se inscreveu.
      \end{definicao}
      \scriptsize
      \begin{tabularx}{\linewidth}{@{}l L l l@{}}
        \toprule
        \cabecalho{Tipo} & \cabecalho{Comportamento} & \cabecalho{Kotlin} & \cabecalho{Dart} \\
        \midrule
        Frio & o produtor começa a cada inscrição & \cod{Flow} & \cod{Stream} (\cod{async*}) \\
        Quente & existe sem observadores; quem chega tarde perde o passado & \cod{SharedFlow} & \cod{Stream} \emph{broadcast} \\
        Estado & quente e guarda o valor atual & \cod{StateFlow} & \cod{ValueNotifier} \\
        \bottomrule
      \end{tabularx}
    \end{column}
    \begin{column}{0.43\textwidth}
      \begin{definicao}[Contrapressão (\emph{backpressure})]
        O que fazer quando o produtor é mais rápido que o consumidor: acumular (\emph{buffer}), descartar, manter só o mais recente (\emph{conflate}) ou amostrar (\emph{debounce}, \emph{sample}).
      \end{definicao}
      \begin{exemplo}[No LabMonitor]
        Um sensor envia 100 leituras por segundo e a tela desenha 60 quadros: basta mostrar sempre a leitura mais recente.
      \end{exemplo}
    \end{column}
  \end{columns}
\end{frame}

% =============================================================================
\section{Fundamentos: concorrência e assincronia}

\begin{frame}{Programas, processos e \emph{threads}}
  \begin{columns}[c]
    \begin{column}{0.47\textwidth}
      \centering
      \begin{tikzpicture}[every node/.style={font=\tiny}]
        \draw[azulpetroleo,line width=0.8pt,rounded corners=4pt,fill=cinzafundo] (0,0) rectangle (5.6,3.2);
        \node[rotulo,anchor=north west] at (0.1,3.15) {\textbf{processo do app}};
        \node[caixaambar,minimum width=5.1cm,minimum height=5mm,font=\tiny] at (2.8,2.45) {memória compartilhada (\emph{heap}): objetos do app};
        \foreach \x/\n in {0.95/{principal (UI)},2.8/{rede},4.65/{cálculo}}{
          \node[caixaescura,minimum width=1.6cm,minimum height=11mm,font=\tiny,align=center] at (\x,1.15) {\emph{thread}\\ \n\\[1pt] \textcolor{ambarclaro}{pilha própria}};
        }
        \node[caixa,minimum width=5.6cm,minimum height=5mm,font=\tiny] (so) at (2.8,-0.55) {sistema operacional: escalonador};
        \foreach \x in {0.7,2.1,3.5,4.9}{\node[caixa,minimum width=1.1cm,minimum height=4mm,font=\tiny] at (\x,-1.3) {núcleo};}
      \end{tikzpicture}
    \end{column}
    \begin{column}{0.51\textwidth}
      \begin{definicao}[Processo e \emph{thread}]
        \textbf{Processo}: programa em execução, com memória própria e isolada dos outros apps. \textbf{\emph{Thread}}: linha de execução dentro do processo; as \emph{threads} dividem a memória do processo, mas cada uma tem sua pilha.
      \end{definicao}
      \begin{pordentro}[O escalonador]
        O sistema operacional distribui as \emph{threads} pelos núcleos e alterna entre elas milhares de vezes por segundo (troca de contexto). Cada \emph{thread} custa memória e tempo de troca: não se cria uma por tarefa.
      \end{pordentro}
      \begin{analogia}
        O processo é a cozinha; as \emph{threads} são os cozinheiros. Eles dividem a despensa, e é aí que começam as brigas pelo mesmo ingrediente.
      \end{analogia}
    \end{column}
  \end{columns}
\end{frame}

\begin{frame}{Concorrência e paralelismo}
  \begin{columns}[c]
    \begin{column}{0.47\textwidth}
      \centering
      \begin{tikzpicture}[x=0.55cm,y=0.5cm,every node/.style={font=\tiny}]
        \node[rotulo,anchor=east] at (-0.2,2) {1 núcleo};
        \foreach \x/\cor in {0/petroleo,1.5/ambar,3/petroleo,4.5/ambar,6/petroleo,7.5/ambar}{\fill[\cor!80] (\x,1.6) rectangle ++(1.4,0.8);}
        \node[rotulo] at (4.5,3.1) {\textbf{concorrência}: as tarefas se alternam};
        \node[rotulo,anchor=east] at (-0.2,0) {núcleo 1};
        \fill[petroleo!80] (0,-0.4) rectangle ++(8.9,0.8);
        \node[rotulo,anchor=east] at (-0.2,-1.1) {núcleo 2};
        \fill[ambar!80] (0,-1.5) rectangle ++(8.9,0.8);
        \node[rotulo] at (4.5,-2.3) {\textbf{paralelismo}: as tarefas rodam juntas};
        \draw[->,cinzamedio] (0,-2.9) -- (9.3,-2.9) node[rotulo,right]{tempo};
        \fill[petroleo!80] (0.4,4) rectangle ++(0.5,0.4); \node[rotulo,anchor=west] at (1,4.2) {tarefa A};
        \fill[ambar!80] (3.4,4) rectangle ++(0.5,0.4); \node[rotulo,anchor=west] at (4,4.2) {tarefa B};
      \end{tikzpicture}
      \raggedright
      \begin{pordentro}[Chips de celular]
        Combinam núcleos de alto desempenho e núcleos econômicos (big.LITTLE): o sistema escolhe onde cada \emph{thread} roda para poupar bateria.
      \end{pordentro}
    \end{column}
    \begin{column}{0.51\textwidth}
      \begin{definicao}[Concorrência]
        Várias tarefas em andamento no mesmo período, que podem se alternar. É um jeito de \textbf{estruturar} o programa.
      \end{definicao}
      \begin{definicao}[Paralelismo]
        Várias tarefas executando \textbf{no mesmo instante}, em núcleos diferentes. É um jeito de \textbf{executar}.
      \end{definicao}
      \begin{historia}[Na frase de Rob Pike (2012)]
        Concorrência é lidar com muitas coisas ao mesmo tempo; paralelismo é fazer muitas coisas ao mesmo tempo.
      \end{historia}
    \end{column}
  \end{columns}
\end{frame}

\begin{frame}{Tarefas de I/O e tarefas de CPU}
  \scriptsize
  \begin{tabularx}{\textwidth}{@{}l L L l l@{}}
    \toprule
    \cabecalho{Tipo} & \cabecalho{O que limita} & \cabecalho{Exemplos} & \cabecalho{Kotlin} & \cabecalho{Dart} \\
    \midrule
    I/O & espera por rede, disco ou banco & requisição HTTP, consulta SQL, ler arquivo & \cod{suspend} + \cod{Dispatchers.IO} & \cod{Future} + \cod{await} \\
    CPU & cálculo no processador & decodificar imagem, JSON gigante, criptografia & \cod{Dispatchers.Default} & \cod{Isolate.run} \\
    Tempo & o relógio & atraso, repetição, tempo limite & \cod{delay}, \cod{withTimeout} & \cod{Future.delayed}, \cod{Timer} \\
    \bottomrule
  \end{tabularx}
  \normalsize
  \vspace{4pt}
  \begin{columns}[T]
    \begin{column}{0.49\textwidth}
      \begin{definicao}[Bloqueante \emph{versus} não bloqueante]
        Uma chamada \textbf{bloqueante} prende a \emph{thread} até terminar. Uma \textbf{não bloqueante} registra o pedido e libera a \emph{thread}; o resultado chega depois, por notificação do sistema.
      \end{definicao}
    \end{column}
    \begin{column}{0.49\textwidth}
      \begin{conceito}[A ideia-chave do módulo]
        \textbf{Esperar não precisa de \emph{thread}}: enquanto a resposta não chega, a mesma \emph{thread} atende outros eventos. Já o \textbf{cálculo pesado precisa de um núcleo} de verdade, fora da \emph{thread} da interface.
      \end{conceito}
    \end{column}
  \end{columns}
\end{frame}

\begin{frame}{O laço de eventos em cada plataforma}
  \begin{columns}[T]
    \begin{column}{0.55\textwidth}
      \begin{definicao}[Laço de eventos (\emph{event loop})]
        Uma \emph{thread} que repete sem parar: pega o próximo evento da fila, executa o código associado até o fim e volta à fila. Toda interface gráfica funciona assim.
      \end{definicao}
      \scriptsize
      \begin{tabularx}{\linewidth}{@{}l L@{}}
        \toprule
        \cabecalho{Plataforma} & \cabecalho{Mecanismo na \emph{thread} principal} \\
        \midrule
        Android & \cod{Looper} + \cod{MessageQueue} + \cod{Handler} \\
        iOS & \cod{RunLoop} e fila principal do GCD \\
        Dart & fila de microtarefas e fila de eventos \\
        Web & laço de eventos do navegador (tarefas e microtarefas) \\
        \bottomrule
      \end{tabularx}
    \end{column}
    \begin{column}{0.43\textwidth}
      \begin{pordentro}[Dart: duas filas]
        Antes de cada evento, o Dart esvazia a fila de \textbf{microtarefas} (continuações curtas de \cod{Future}s). Desde o Flutter 3.29, no Android e no iOS, o código Dart roda na própria \emph{thread} principal do app.
      \end{pordentro}
      \begin{atencao}[Um evento lento atrasa todos]
        Quadros se perdem (\emph{jank}) e, no Android, sem resposta a um toque por cerca de 5\,s aparece o diálogo ANR.
      \end{atencao}
    \end{column}
  \end{columns}
\end{frame}

\begin{frame}{Da pirâmide de \emph{callbacks} ao \cod{async}/\cod{await}}
  \centering
  \begin{tikzpicture}[every node/.style={font=\tiny}]
    \draw[seta] (0,0) -- (13.4,0);
    \foreach \x/\t/\d in {1.2/\emph{callbacks}/{funções chamadas\\ no fim da tarefa},4.4/\emph{Futures} e \emph{Promises}/{objeto que representa\\ o resultado futuro},7.8/{\cod{async}/\cod{await}}/{C\# (2012), Dart (2015),\\ JavaScript (2017)},11.4/concorrência estruturada/{Kotlin (2018):\\ escopos e cancelamento}}{
      \fill[ambar] (\x,0) circle (2pt);
      \node[font=\scriptsize\bfseries,text=azulpetroleo,above=1mm] at (\x,0) {\t};
      \node[rotulo,below=1mm] at (\x,0) {\d};
    }
  \end{tikzpicture}
  \raggedright
  \begin{columns}[T]
    \begin{column}{0.49\textwidth}
      \begin{definicao}[\emph{Future} (ou \emph{Promise})]
        Objeto que representa um resultado que ainda não existe. Começa \textbf{pendente} e termina \textbf{com valor} ou \textbf{com erro}. Quem o recebe pode esperar sem bloquear a \emph{thread}.
      \end{definicao}
      \centering
      \begin{tikzpicture}[x=0.5cm,y=0.45cm,every node/.style={font=\tiny}]
        \node[rotulo,anchor=east] at (-0.2,0.4) {\emph{thread}};
        \fill[petroleo!80] (0,0) rectangle ++(2,0.8); \node[white] at (1,0.4) {trecho 1};
        \fill[cinzaclaro] (2,0) rectangle ++(3,0.8); \node at (3.5,0.4) {outros eventos};
        \fill[petroleo!80] (5,0) rectangle ++(2,0.8); \node[white] at (6,0.4) {trecho 2};
        \fill[cinzaclaro] (7,0) rectangle ++(2.4,0.8); \node at (8.2,0.4) {outros};
        \fill[petroleo!80] (9.4,0) rectangle ++(1.6,0.8); \node[white] at (10.2,0.4) {fim};
        \foreach \x in {2,7}{\draw[ambar,line width=1pt] (\x,-0.2) -- (\x,1.1); \node[rotulo,above] at (\x,1.1) {\cod{await}};}
      \end{tikzpicture}
    \end{column}
    \begin{column}{0.49\textwidth}
      \begin{pordentro}[O que acontece em um \cod{await}]
        A função pausa em um \textbf{ponto de suspensão} e devolve a \emph{thread}; quando o resultado chega, continua de onde parou, com as variáveis locais preservadas. O compilador do Kotlin transforma cada função \cod{suspend} em uma máquina de estados que recebe uma \cod{Continuation}.
      \end{pordentro}
      \begin{dica}
        O ganho é escrever código assíncrono com cara de sequencial: \cod{try}/\cod{catch} e laços funcionam como sempre.
      \end{dica}
    \end{column}
  \end{columns}
\end{frame}

\begin{frame}{Concorrência estruturada}
  \begin{columns}[c]
    \begin{column}{0.45\textwidth}
      \centering
      \begin{tikzpicture}[every node/.style={font=\tiny}]
        \node[caixaescura,minimum width=3.4cm,font=\tiny] (e) at (0,0) {escopo da tela};
        \node[caixa,minimum width=1.6cm,font=\tiny] (a) at (-0.95,-1) {carregar lista};
        \node[caixa,minimum width=1.6cm,font=\tiny] (b) at (0.95,-1) {ouvir sensor};
        \node[caixaambar,minimum width=1.6cm,font=\tiny] (c) at (0.95,-2) {ler histórico};
        \draw[seta] (e) -- (a); \draw[seta] (e) -- (b); \draw[seta] (b) -- (c);
        \draw[pontilhado] (-2,0.5) rectangle (2,-2.5);
        \node[rotulo,text width=4.2cm] at (0,-3) {a tela fecha $\Rightarrow$ o escopo termina $\Rightarrow$ as tarefas filhas são canceladas};
      \end{tikzpicture}
      \begin{historia}
        Assim como a programação estruturada aposentou o \cod{goto}, este modelo aposenta tarefas soltas. A ideia ganhou força em 2018, com um ensaio de Nathaniel J. Smith, e o Kotlin a adotou no mesmo ano.
      \end{historia}
    \end{column}
    \begin{column}{0.53\textwidth}
      \begin{definicao}[Concorrência estruturada]
        Toda tarefa concorrente nasce dentro de um \textbf{escopo} e não vive mais que ele: quando o escopo termina, as tarefas filhas são canceladas.
      \end{definicao}
      \begin{conceito}[Garantias]
        Nada fica rodando depois que a tela fecha; o cancelamento chega aos filhos; erros dos filhos chegam ao pai.
      \end{conceito}
      \begin{atencao}[E no Dart?]
        \cod{Future} e \cod{Stream} não seguem esse modelo: cancele inscrições e temporizadores no \cod{dispose} e confira \cod{mounted} após um \cod{await}.
      \end{atencao}
    \end{column}
  \end{columns}
\end{frame}

\begin{frame}{Condições de corrida e impasses}
  \begin{columns}[T]
    \begin{column}{0.47\textwidth}
      \centering
      \begin{tikzpicture}[x=0.7cm,y=0.6cm,every node/.style={font=\tiny}]
        \tikzset{passo/.style={caixa,minimum height=4.5mm,minimum width=1.2cm,inner sep=2pt,font=\tiny}}
        \node[rotulo,anchor=east] at (-0.1,1) {\emph{thread} A};
        \node[rotulo,anchor=east] at (-0.1,0) {\emph{thread} B};
        \node[passo] at (0.9,1) {lê 5};
        \node[passo] at (2.6,0) {lê 5};
        \node[passo,caixaambar,font=\tiny] at (4.3,1) {grava 6};
        \node[passo,caixaambar,font=\tiny] at (6,0) {grava 6};
        \draw[->,cinzamedio] (0,-0.7) -- (7,-0.7) node[rotulo,right]{t};
        \node[rotulo,text=vermelhoerro,text width=4.8cm] at (3.4,-1.3) {duas somas de 1 a partir de 5 deveriam dar 7: uma atualização se perdeu};
      \end{tikzpicture}
      \begin{definicao}[Condição de corrida]
        Erro em que o resultado depende da ordem imprevisível em que tarefas concorrentes acessam um dado compartilhado.
      \end{definicao}
      \begin{atencao}[Impasse (\emph{deadlock})]
        Duas tarefas, cada uma segurando um recurso e esperando o da outra: ninguém avança. Adquira travas sempre na mesma ordem.
      \end{atencao}
    \end{column}
    \begin{column}{0.51\textwidth}
      \scriptsize
      \begin{tabularx}{\linewidth}{@{}l L@{}}
        \toprule
        \cabecalho{Proteção} & \cabecalho{Como funciona} \\
        \midrule
        Exclusão mútua & uma tarefa por vez na região crítica (\cod{Mutex}) \\
        Operação atômica & ler e gravar em um passo indivisível (\cod{AtomicInteger}, \cod{StateFlow.update}) \\
        Confinamento & só uma \emph{thread} toca o dado (a principal, por exemplo) \\
        Imutabilidade & ninguém altera: cria-se uma cópia nova \\
        Mensagens & cada tarefa tem seus dados e troca cópias \\
        \bottomrule
      \end{tabularx}
      \normalsize
      \begin{pordentro}[Corrida lógica com uma \emph{thread} só]
        Entre dois \cod{await}, outro evento pode mudar o estado. Duas buscas disparadas em sequência podem responder fora de ordem, e a mais antiga sobrescrever a mais nova.
      \end{pordentro}
    \end{column}
  \end{columns}
\end{frame}

\begin{frame}{Memória compartilhada ou troca de mensagens}
  \begin{columns}[c]
    \begin{column}{0.43\textwidth}
      \centering
      \begin{tikzpicture}[every node/.style={font=\tiny}]
        \node[caixa,minimum width=1.3cm,font=\tiny] (a) at (0,0) {\emph{thread} A};
        \node[caixa,minimum width=1.3cm,font=\tiny] (b) at (2.6,0) {\emph{thread} B};
        \node[caixaambar,minimum width=1.6cm,font=\tiny] (o) at (1.3,-1.2) {objeto + trava};
        \draw[seta] (a) -- (o); \draw[seta] (b) -- (o);
        \node[rotulo] at (1.3,0.6) {\textbf{memória compartilhada}};
        \node[caixaescura,minimum width=1.3cm,minimum height=9mm,font=\tiny] (c) at (0,-3) {isolate A\\ memória A};
        \node[caixaescura,minimum width=1.3cm,minimum height=9mm,font=\tiny] (d) at (2.6,-3) {isolate B\\ memória B};
        \draw[setaambar] ([yshift=1.5mm]c.east) -- node[rotulo,above]{mensagem} ([yshift=1.5mm]d.west);
        \draw[setaambar] ([yshift=-1.5mm]d.west) -- node[rotulo,below]{resposta} ([yshift=-1.5mm]c.east);
        \node[rotulo] at (1.3,-2.1) {\textbf{troca de mensagens}};
      \end{tikzpicture}
    \end{column}
    \begin{column}{0.55\textwidth}
      \scriptsize
      \begin{tabularx}{\linewidth}{@{}l L L@{}}
        \toprule
         & \cabecalho{Memória compartilhada} & \cabecalho{Troca de mensagens} \\
        \midrule
        Exemplos & \emph{threads} da JVM (Kotlin, Java) & \emph{isolates} do Dart, atores, Web Workers \\
        Comunicação & ler e gravar o mesmo objeto & enviar cópias de dados \\
        Risco & corridas e impasses & custo de copiar dados grandes \\
        Proteção & travas, atômicos, confinamento & o próprio modelo isola \\
        \bottomrule
      \end{tabularx}
      \normalsize
      \begin{analogia}
        Um quadro branco compartilhado exige regras para ninguém apagar o do colega. Colegas que trocam bilhetes nunca brigam pelo quadro, mas gastam tempo escrevendo.
      \end{analogia}
    \end{column}
  \end{columns}
\end{frame}

\begin{frame}{Fluxos de dados: frios, quentes e contrapressão}
  \begin{columns}[T]
    \begin{column}{0.55\textwidth}
      \begin{definicao}[Fluxo (\emph{stream})]
        Sequência de valores que chegam ao longo do tempo: leituras de sensor, mensagens, mudanças no banco. Segue o padrão \textbf{Observador}: quem produz notifica quem se inscreveu.
      \end{definicao}
      \scriptsize
      \begin{tabularx}{\linewidth}{@{}l L l l@{}}
        \toprule
        \cabecalho{Tipo} & \cabecalho{Comportamento} & \cabecalho{Kotlin} & \cabecalho{Dart} \\
        \midrule
        Frio & o produtor começa a cada inscrição & \cod{Flow} & \cod{Stream} (\cod{async*}) \\
        Quente & existe sem observadores; quem chega tarde perde o passado & \cod{SharedFlow} & \cod{Stream} \emph{broadcast} \\
        Estado & quente e guarda o valor atual & \cod{StateFlow} & \cod{ValueNotifier} \\
        \bottomrule
      \end{tabularx}
    \end{column}
    \begin{column}{0.43\textwidth}
      \begin{definicao}[Contrapressão (\emph{backpressure})]
        O que fazer quando o produtor é mais rápido que o consumidor: acumular (\emph{buffer}), descartar, manter só o mais recente (\emph{conflate}) ou amostrar (\emph{debounce}, \emph{sample}).
      \end{definicao}
      \begin{exemplo}[No LabMonitor]
        Um sensor envia 100 leituras por segundo e a tela desenha 60 quadros: basta mostrar sempre a leitura mais recente.
      \end{exemplo}
    \end{column}
  \end{columns}
\end{frame}

% =============================================================================
\section{Fundamentos: concorrência e assincronia}

\begin{frame}{Programas, processos e \emph{threads}}
  \begin{columns}[c]
    \begin{column}{0.47\textwidth}
      \centering
      \begin{tikzpicture}[every node/.style={font=\tiny}]
        \draw[azulpetroleo,line width=0.8pt,rounded corners=4pt,fill=cinzafundo] (0,0) rectangle (5.6,3.2);
        \node[rotulo,anchor=north west] at (0.1,3.15) {\textbf{processo do app}};
        \node[caixaambar,minimum width=5.1cm,minimum height=5mm,font=\tiny] at (2.8,2.45) {memória compartilhada (\emph{heap}): objetos do app};
        \foreach \x/\n in {0.95/{principal (UI)},2.8/{rede},4.65/{cálculo}}{
          \node[caixaescura,minimum width=1.6cm,minimum height=11mm,font=\tiny,align=center] at (\x,1.15) {\emph{thread}\\ \n\\[1pt] \textcolor{ambarclaro}{pilha própria}};
        }
        \node[caixa,minimum width=5.6cm,minimum height=5mm,font=\tiny] (so) at (2.8,-0.55) {sistema operacional: escalonador};
        \foreach \x in {0.7,2.1,3.5,4.9}{\node[caixa,minimum width=1.1cm,minimum height=4mm,font=\tiny] at (\x,-1.3) {núcleo};}
      \end{tikzpicture}
    \end{column}
    \begin{column}{0.51\textwidth}
      \begin{definicao}[Processo e \emph{thread}]
        \textbf{Processo}: programa em execução, com memória própria e isolada dos outros apps. \textbf{\emph{Thread}}: linha de execução dentro do processo; as \emph{threads} dividem a memória do processo, mas cada uma tem sua pilha.
      \end{definicao}
      \begin{pordentro}[O escalonador]
        O sistema operacional distribui as \emph{threads} pelos núcleos e alterna entre elas milhares de vezes por segundo (troca de contexto). Cada \emph{thread} custa memória e tempo de troca: não se cria uma por tarefa.
      \end{pordentro}
      \begin{analogia}
        O processo é a cozinha; as \emph{threads} são os cozinheiros. Eles dividem a despensa, e é aí que começam as brigas pelo mesmo ingrediente.
      \end{analogia}
    \end{column}
  \end{columns}
\end{frame}

\begin{frame}{Concorrência e paralelismo}
  \begin{columns}[c]
    \begin{column}{0.47\textwidth}
      \centering
      \begin{tikzpicture}[x=0.55cm,y=0.5cm,every node/.style={font=\tiny}]
        \node[rotulo,anchor=east] at (-0.2,2) {1 núcleo};
        \foreach \x/\cor in {0/petroleo,1.5/ambar,3/petroleo,4.5/ambar,6/petroleo,7.5/ambar}{\fill[\cor!80] (\x,1.6) rectangle ++(1.4,0.8);}
        \node[rotulo] at (4.5,3.1) {\textbf{concorrência}: as tarefas se alternam};
        \node[rotulo,anchor=east] at (-0.2,0) {núcleo 1};
        \fill[petroleo!80] (0,-0.4) rectangle ++(8.9,0.8);
        \node[rotulo,anchor=east] at (-0.2,-1.1) {núcleo 2};
        \fill[ambar!80] (0,-1.5) rectangle ++(8.9,0.8);
        \node[rotulo] at (4.5,-2.3) {\textbf{paralelismo}: as tarefas rodam juntas};
        \draw[->,cinzamedio] (0,-2.9) -- (9.3,-2.9) node[rotulo,right]{tempo};
        \fill[petroleo!80] (0.4,4) rectangle ++(0.5,0.4); \node[rotulo,anchor=west] at (1,4.2) {tarefa A};
        \fill[ambar!80] (3.4,4) rectangle ++(0.5,0.4); \node[rotulo,anchor=west] at (4,4.2) {tarefa B};
      \end{tikzpicture}
      \raggedright
      \begin{pordentro}[Chips de celular]
        Combinam núcleos de alto desempenho e núcleos econômicos (big.LITTLE): o sistema escolhe onde cada \emph{thread} roda para poupar bateria.
      \end{pordentro}
    \end{column}
    \begin{column}{0.51\textwidth}
      \begin{definicao}[Concorrência]
        Várias tarefas em andamento no mesmo período, que podem se alternar. É um jeito de \textbf{estruturar} o programa.
      \end{definicao}
      \begin{definicao}[Paralelismo]
        Várias tarefas executando \textbf{no mesmo instante}, em núcleos diferentes. É um jeito de \textbf{executar}.
      \end{definicao}
      \begin{historia}[Na frase de Rob Pike (2012)]
        Concorrência é lidar com muitas coisas ao mesmo tempo; paralelismo é fazer muitas coisas ao mesmo tempo.
      \end{historia}
    \end{column}
  \end{columns}
\end{frame}

\begin{frame}{Tarefas de I/O e tarefas de CPU}
  \scriptsize
  \begin{tabularx}{\textwidth}{@{}l L L l l@{}}
    \toprule
    \cabecalho{Tipo} & \cabecalho{O que limita} & \cabecalho{Exemplos} & \cabecalho{Kotlin} & \cabecalho{Dart} \\
    \midrule
    I/O & espera por rede, disco ou banco & requisição HTTP, consulta SQL, ler arquivo & \cod{suspend} + \cod{Dispatchers.IO} & \cod{Future} + \cod{await} \\
    CPU & cálculo no processador & decodificar imagem, JSON gigante, criptografia & \cod{Dispatchers.Default} & \cod{Isolate.run} \\
    Tempo & o relógio & atraso, repetição, tempo limite & \cod{delay}, \cod{withTimeout} & \cod{Future.delayed}, \cod{Timer} \\
    \bottomrule
  \end{tabularx}
  \normalsize
  \vspace{4pt}
  \begin{columns}[T]
    \begin{column}{0.49\textwidth}
      \begin{definicao}[Bloqueante \emph{versus} não bloqueante]
        Uma chamada \textbf{bloqueante} prende a \emph{thread} até terminar. Uma \textbf{não bloqueante} registra o pedido e libera a \emph{thread}; o resultado chega depois, por notificação do sistema.
      \end{definicao}
    \end{column}
    \begin{column}{0.49\textwidth}
      \begin{conceito}[A ideia-chave do módulo]
        \textbf{Esperar não precisa de \emph{thread}}: enquanto a resposta não chega, a mesma \emph{thread} atende outros eventos. Já o \textbf{cálculo pesado precisa de um núcleo} de verdade, fora da \emph{thread} da interface.
      \end{conceito}
    \end{column}
  \end{columns}
\end{frame}

\begin{frame}{O laço de eventos em cada plataforma}
  \begin{columns}[T]
    \begin{column}{0.55\textwidth}
      \begin{definicao}[Laço de eventos (\emph{event loop})]
        Uma \emph{thread} que repete sem parar: pega o próximo evento da fila, executa o código associado até o fim e volta à fila. Toda interface gráfica funciona assim.
      \end{definicao}
      \scriptsize
      \begin{tabularx}{\linewidth}{@{}l L@{}}
        \toprule
        \cabecalho{Plataforma} & \cabecalho{Mecanismo na \emph{thread} principal} \\
        \midrule
        Android & \cod{Looper} + \cod{MessageQueue} + \cod{Handler} \\
        iOS & \cod{RunLoop} e fila principal do GCD \\
        Dart & fila de microtarefas e fila de eventos \\
        Web & laço de eventos do navegador (tarefas e microtarefas) \\
        \bottomrule
      \end{tabularx}
    \end{column}
    \begin{column}{0.43\textwidth}
      \begin{pordentro}[Dart: duas filas]
        Antes de cada evento, o Dart esvazia a fila de \textbf{microtarefas} (continuações curtas de \cod{Future}s). Desde o Flutter 3.29, no Android e no iOS, o código Dart roda na própria \emph{thread} principal do app.
      \end{pordentro}
      \begin{atencao}[Um evento lento atrasa todos]
        Quadros se perdem (\emph{jank}) e, no Android, sem resposta a um toque por cerca de 5\,s aparece o diálogo ANR.
      \end{atencao}
    \end{column}
  \end{columns}
\end{frame}

\begin{frame}{Da pirâmide de \emph{callbacks} ao \cod{async}/\cod{await}}
  \centering
  \begin{tikzpicture}[every node/.style={font=\tiny}]
    \draw[seta] (0,0) -- (13.4,0);
    \foreach \x/\t/\d in {1.2/\emph{callbacks}/{funções chamadas\\ no fim da tarefa},4.4/\emph{Futures} e \emph{Promises}/{objeto que representa\\ o resultado futuro},7.8/{\cod{async}/\cod{await}}/{C\# (2012), Dart (2015),\\ JavaScript (2017)},11.4/concorrência estruturada/{Kotlin (2018):\\ escopos e cancelamento}}{
      \fill[ambar] (\x,0) circle (2pt);
      \node[font=\scriptsize\bfseries,text=azulpetroleo,above=1mm] at (\x,0) {\t};
      \node[rotulo,below=1mm] at (\x,0) {\d};
    }
  \end{tikzpicture}
  \raggedright
  \begin{columns}[T]
    \begin{column}{0.49\textwidth}
      \begin{definicao}[\emph{Future} (ou \emph{Promise})]
        Objeto que representa um resultado que ainda não existe. Começa \textbf{pendente} e termina \textbf{com valor} ou \textbf{com erro}. Quem o recebe pode esperar sem bloquear a \emph{thread}.
      \end{definicao}
      \centering
      \begin{tikzpicture}[x=0.5cm,y=0.45cm,every node/.style={font=\tiny}]
        \node[rotulo,anchor=east] at (-0.2,0.4) {\emph{thread}};
        \fill[petroleo!80] (0,0) rectangle ++(2,0.8); \node[white] at (1,0.4) {trecho 1};
        \fill[cinzaclaro] (2,0) rectangle ++(3,0.8); \node at (3.5,0.4) {outros eventos};
        \fill[petroleo!80] (5,0) rectangle ++(2,0.8); \node[white] at (6,0.4) {trecho 2};
        \fill[cinzaclaro] (7,0) rectangle ++(2.4,0.8); \node at (8.2,0.4) {outros};
        \fill[petroleo!80] (9.4,0) rectangle ++(1.6,0.8); \node[white] at (10.2,0.4) {fim};
        \foreach \x in {2,7}{\draw[ambar,line width=1pt] (\x,-0.2) -- (\x,1.1); \node[rotulo,above] at (\x,1.1) {\cod{await}};}
      \end{tikzpicture}
    \end{column}
    \begin{column}{0.49\textwidth}
      \begin{pordentro}[O que acontece em um \cod{await}]
        A função pausa em um \textbf{ponto de suspensão} e devolve a \emph{thread}; quando o resultado chega, continua de onde parou, com as variáveis locais preservadas. O compilador do Kotlin transforma cada função \cod{suspend} em uma máquina de estados que recebe uma \cod{Continuation}.
      \end{pordentro}
      \begin{dica}
        O ganho é escrever código assíncrono com cara de sequencial: \cod{try}/\cod{catch} e laços funcionam como sempre.
      \end{dica}
    \end{column}
  \end{columns}
\end{frame}

\begin{frame}{Concorrência estruturada}
  \begin{columns}[c]
    \begin{column}{0.45\textwidth}
      \centering
      \begin{tikzpicture}[every node/.style={font=\tiny}]
        \node[caixaescura,minimum width=3.4cm,font=\tiny] (e) at (0,0) {escopo da tela};
        \node[caixa,minimum width=1.6cm,font=\tiny] (a) at (-0.95,-1) {carregar lista};
        \node[caixa,minimum width=1.6cm,font=\tiny] (b) at (0.95,-1) {ouvir sensor};
        \node[caixaambar,minimum width=1.6cm,font=\tiny] (c) at (0.95,-2) {ler histórico};
        \draw[seta] (e) -- (a); \draw[seta] (e) -- (b); \draw[seta] (b) -- (c);
        \draw[pontilhado] (-2,0.5) rectangle (2,-2.5);
        \node[rotulo,text width=4.2cm] at (0,-3) {a tela fecha $\Rightarrow$ o escopo termina $\Rightarrow$ as tarefas filhas são canceladas};
      \end{tikzpicture}
      \begin{historia}
        Assim como a programação estruturada aposentou o \cod{goto}, este modelo aposenta tarefas soltas. A ideia ganhou força em 2018, com um ensaio de Nathaniel J. Smith, e o Kotlin a adotou no mesmo ano.
      \end{historia}
    \end{column}
    \begin{column}{0.53\textwidth}
      \begin{definicao}[Concorrência estruturada]
        Toda tarefa concorrente nasce dentro de um \textbf{escopo} e não vive mais que ele: quando o escopo termina, as tarefas filhas são canceladas.
      \end{definicao}
      \begin{conceito}[Garantias]
        Nada fica rodando depois que a tela fecha; o cancelamento chega aos filhos; erros dos filhos chegam ao pai.
      \end{conceito}
      \begin{atencao}[E no Dart?]
        \cod{Future} e \cod{Stream} não seguem esse modelo: cancele inscrições e temporizadores no \cod{dispose} e confira \cod{mounted} após um \cod{await}.
      \end{atencao}
    \end{column}
  \end{columns}
\end{frame}

\begin{frame}{Condições de corrida e impasses}
  \begin{columns}[T]
    \begin{column}{0.47\textwidth}
      \centering
      \begin{tikzpicture}[x=0.7cm,y=0.6cm,every node/.style={font=\tiny}]
        \tikzset{passo/.style={caixa,minimum height=4.5mm,minimum width=1.2cm,inner sep=2pt,font=\tiny}}
        \node[rotulo,anchor=east] at (-0.1,1) {\emph{thread} A};
        \node[rotulo,anchor=east] at (-0.1,0) {\emph{thread} B};
        \node[passo] at (0.9,1) {lê 5};
        \node[passo] at (2.6,0) {lê 5};
        \node[passo,caixaambar,font=\tiny] at (4.3,1) {grava 6};
        \node[passo,caixaambar,font=\tiny] at (6,0) {grava 6};
        \draw[->,cinzamedio] (0,-0.7) -- (7,-0.7) node[rotulo,right]{t};
        \node[rotulo,text=vermelhoerro,text width=4.8cm] at (3.4,-1.3) {duas somas de 1 a partir de 5 deveriam dar 7: uma atualização se perdeu};
      \end{tikzpicture}
      \begin{definicao}[Condição de corrida]
        Erro em que o resultado depende da ordem imprevisível em que tarefas concorrentes acessam um dado compartilhado.
      \end{definicao}
      \begin{atencao}[Impasse (\emph{deadlock})]
        Duas tarefas, cada uma segurando um recurso e esperando o da outra: ninguém avança. Adquira travas sempre na mesma ordem.
      \end{atencao}
    \end{column}
    \begin{column}{0.51\textwidth}
      \scriptsize
      \begin{tabularx}{\linewidth}{@{}l L@{}}
        \toprule
        \cabecalho{Proteção} & \cabecalho{Como funciona} \\
        \midrule
        Exclusão mútua & uma tarefa por vez na região crítica (\cod{Mutex}) \\
        Operação atômica & ler e gravar em um passo indivisível (\cod{AtomicInteger}, \cod{StateFlow.update}) \\
        Confinamento & só uma \emph{thread} toca o dado (a principal, por exemplo) \\
        Imutabilidade & ninguém altera: cria-se uma cópia nova \\
        Mensagens & cada tarefa tem seus dados e troca cópias \\
        \bottomrule
      \end{tabularx}
      \normalsize
      \begin{pordentro}[Corrida lógica com uma \emph{thread} só]
        Entre dois \cod{await}, outro evento pode mudar o estado. Duas buscas disparadas em sequência podem responder fora de ordem, e a mais antiga sobrescrever a mais nova.
      \end{pordentro}
    \end{column}
  \end{columns}
\end{frame}

\begin{frame}{Memória compartilhada ou troca de mensagens}
  \begin{columns}[c]
    \begin{column}{0.43\textwidth}
      \centering
      \begin{tikzpicture}[every node/.style={font=\tiny}]
        \node[caixa,minimum width=1.3cm,font=\tiny] (a) at (0,0) {\emph{thread} A};
        \node[caixa,minimum width=1.3cm,font=\tiny] (b) at (2.6,0) {\emph{thread} B};
        \node[caixaambar,minimum width=1.6cm,font=\tiny] (o) at (1.3,-1.2) {objeto + trava};
        \draw[seta] (a) -- (o); \draw[seta] (b) -- (o);
        \node[rotulo] at (1.3,0.6) {\textbf{memória compartilhada}};
        \node[caixaescura,minimum width=1.3cm,minimum height=9mm,font=\tiny] (c) at (0,-3) {isolate A\\ memória A};
        \node[caixaescura,minimum width=1.3cm,minimum height=9mm,font=\tiny] (d) at (2.6,-3) {isolate B\\ memória B};
        \draw[setaambar] ([yshift=1.5mm]c.east) -- node[rotulo,above]{mensagem} ([yshift=1.5mm]d.west);
        \draw[setaambar] ([yshift=-1.5mm]d.west) -- node[rotulo,below]{resposta} ([yshift=-1.5mm]c.east);
        \node[rotulo] at (1.3,-2.1) {\textbf{troca de mensagens}};
      \end{tikzpicture}
    \end{column}
    \begin{column}{0.55\textwidth}
      \scriptsize
      \begin{tabularx}{\linewidth}{@{}l L L@{}}
        \toprule
         & \cabecalho{Memória compartilhada} & \cabecalho{Troca de mensagens} \\
        \midrule
        Exemplos & \emph{threads} da JVM (Kotlin, Java) & \emph{isolates} do Dart, atores, Web Workers \\
        Comunicação & ler e gravar o mesmo objeto & enviar cópias de dados \\
        Risco & corridas e impasses & custo de copiar dados grandes \\
        Proteção & travas, atômicos, confinamento & o próprio modelo isola \\
        \bottomrule
      \end{tabularx}
      \normalsize
      \begin{analogia}
        Um quadro branco compartilhado exige regras para ninguém apagar o do colega. Colegas que trocam bilhetes nunca brigam pelo quadro, mas gastam tempo escrevendo.
      \end{analogia}
    \end{column}
  \end{columns}
\end{frame}

\begin{frame}{Fluxos de dados: frios, quentes e contrapressão}
  \begin{columns}[T]
    \begin{column}{0.55\textwidth}
      \begin{definicao}[Fluxo (\emph{stream})]
        Sequência de valores que chegam ao longo do tempo: leituras de sensor, mensagens, mudanças no banco. Segue o padrão \textbf{Observador}: quem produz notifica quem se inscreveu.
      \end{definicao}
      \scriptsize
      \begin{tabularx}{\linewidth}{@{}l L l l@{}}
        \toprule
        \cabecalho{Tipo} & \cabecalho{Comportamento} & \cabecalho{Kotlin} & \cabecalho{Dart} \\
        \midrule
        Frio & o produtor começa a cada inscrição & \cod{Flow} & \cod{Stream} (\cod{async*}) \\
        Quente & existe sem observadores; quem chega tarde perde o passado & \cod{SharedFlow} & \cod{Stream} \emph{broadcast} \\
        Estado & quente e guarda o valor atual & \cod{StateFlow} & \cod{ValueNotifier} \\
        \bottomrule
      \end{tabularx}
    \end{column}
    \begin{column}{0.43\textwidth}
      \begin{definicao}[Contrapressão (\emph{backpressure})]
        O que fazer quando o produtor é mais rápido que o consumidor: acumular (\emph{buffer}), descartar, manter só o mais recente (\emph{conflate}) ou amostrar (\emph{debounce}, \emph{sample}).
      \end{definicao}
      \begin{exemplo}[No LabMonitor]
        Um sensor envia 100 leituras por segundo e a tela desenha 60 quadros: basta mostrar sempre a leitura mais recente.
      \end{exemplo}
    \end{column}
  \end{columns}
\end{frame}


% =============================================================================
\section{Por que programação assíncrona?}

\begin{frame}{A \emph{thread} principal (UI thread)}
  \begin{columns}[c]
    \begin{column}{0.52\textwidth}
      \small
      \begin{itemize}
        \item Todo app tem uma \termo{thread principal} que desenha a tela e responde a toques.
        \item Para 60\,Hz, cada quadro precisa ficar pronto em cerca de \textbf{16\,ms}.
        \item Rede, banco de dados, arquivos e processamento pesado levam de dezenas de ms a vários segundos.
        \item Se a thread principal ficar bloqueada por $\sim$5\,s, o sistema mostra o diálogo \termo{ANR} (\emph{Application Not Responding}).
      \end{itemize}
    \end{column}
    \begin{column}{0.46\textwidth}
      \centering
      \begin{tikzpicture}[x=0.52cm,y=0.55cm]
        \node[rotulo,anchor=east] at (0,2) {bloqueada};
        \foreach \x in {0,2}{\fill[petroleo!70] (\x,1.6) rectangle ++(1.8,0.8);}
        \fill[vermelhoerro!80] (4,1.6) rectangle ++(5.5,0.8) node[midway,white,font=\tiny]{rede (2 s)};
        \node[rotulo,anchor=east] at (0,0) {assíncrona};
        \foreach \x in {0,2,4,6,8}{\fill[petroleo!70] (\x,-0.4) rectangle ++(1.8,0.8);}
        \fill[ambar!80] (4,-1.6) rectangle ++(5.5,0.8) node[midway,white,font=\tiny]{rede em outra thread};
        \draw[->,cinzamedio] (0,-2.4) -- (10,-2.4) node[rotulo,right]{t};
        \node[rotulo] at (5,3.2) {cada bloco azul = um quadro desenhado};
      \end{tikzpicture}
    \end{column}
  \end{columns}
  \begin{atencao}[Regra de ouro]
    Nunca faça I/O nem cálculo pesado na thread principal. Desde o Android 3.0, uma requisição de rede na thread principal gera \cod{NetworkOnMainThreadException}.
  \end{atencao}
\end{frame}

\begin{frame}[fragile]{Callbacks \emph{versus} coroutines}
  \begin{columns}[T]
    \begin{column}{0.49\textwidth}
\begin{codigo}[title=Callbacks (\emph{callback hell})]{Kotlin}
login(usuario) { token ->
  buscarBancadas(token) { bancadas ->
    buscarLeituras(bancadas[0]) { leituras ->
      runOnUiThread {
        mostrar(leituras)
      }
    }
  }
}
// erros? cancelamento? cada nível trata o seu
\end{codigo}
    \end{column}
    \begin{column}{0.49\textwidth}
\begin{codigo}[title=Coroutines]{Kotlin}
scope.launch {
  val token = login(usuario)          // suspende
  val bancadas = buscarBancadas(token)
  val leituras = buscarLeituras(bancadas[0])
  mostrar(leituras)                   // de volta à UI
}
// try/catch normal, cancelamento automático
\end{codigo}
    \end{column}
  \end{columns}
  \begin{conceito}[Coroutine]
    Uma computação que pode ser \textbf{suspensa} e \textbf{retomada} sem bloquear a thread. É muito mais leve que uma thread: dá para ter milhares ao mesmo tempo. O código assíncrono fica com cara de código sequencial.
  \end{conceito}
\end{frame}

% =============================================================================
\section{Fundamentos de coroutines}

\begin{frame}[fragile]{Funções \cod{suspend}}
  \begin{columns}[T]
    \begin{column}{0.52\textwidth}
\begin{codigo}[title=Suspensão]{Kotlin}
suspend fun lerTemperatura(): Double {
  delay(500)          // suspende, NÃO bloqueia
  return 24.7
}

suspend fun relatorio(): String {
  val t = lerTemperatura()   // pode chamar suspend
  return "Temperatura: $t"
}

fun naoSuspensa() {
  // lerTemperatura()  // ERRO: só dentro de
  //                   // coroutine ou de suspend
}
\end{codigo}
    \end{column}
    \begin{column}{0.46\textwidth}
      \centering
      \begin{tikzpicture}[node distance=2mm]
        \node[caixaescura,minimum width=4.4cm] (a) {\cod{relatorio()} começa};
        \node[caixaambar,minimum width=4.4cm,below=of a] (b) {\cod{delay(500)}: \textbf{suspende}};
        \node[caixa,minimum width=4.4cm,below=of b,fill=cinzaclaro] (c) {thread livre para outras tarefas};
        \node[caixaambar,minimum width=4.4cm,below=of c] (d) {500 ms depois: \textbf{retoma}};
        \node[caixaescura,minimum width=4.4cm,below=of d] (e) {devolve o texto};
        \draw[seta] (a) -- (b); \draw[seta] (b) -- (c); \draw[seta] (c) -- (d); \draw[seta] (d) -- (e);
      \end{tikzpicture}
      \begin{dica}
        \cod{suspend} não cria thread nem torna nada assíncrono sozinho: só indica que a função \textbf{pode} suspender.
      \end{dica}
    \end{column}
  \end{columns}
\end{frame}

\begin{frame}[fragile]{Construtores: \cod{launch}, \cod{async} e \cod{runBlocking}}
  \begin{columns}[T]
    \begin{column}{0.52\textwidth}
\begin{codigo}[title=Iniciando coroutines]{Kotlin}
// launch: "dispare e esqueça", devolve Job
val job: Job = scope.launch {
  salvarLeitura(leitura)
}
job.cancel()

// async: devolve Deferred<T>, use await()
scope.launch {
  val t = async { lerTemperatura() }
  val u = async { lerUmidade() }
  // as duas rodam em paralelo
  mostrar(t.await(), u.await())
}

// runBlocking: BLOQUEIA a thread.
// Só em main() de exemplos e em testes antigos.
fun main() = runBlocking {
  println(relatorio())
}
\end{codigo}
    \end{column}
    \begin{column}{0.46\textwidth}
      \footnotesize
      \begin{tabularx}{\linewidth}{@{}l l L@{}}
        \toprule
        \cabecalho{Construtor} & \cabecalho{Retorno} & \cabecalho{Uso} \\
        \midrule
        \cod{launch} & \cod{Job} & tarefa sem resultado \\
        \cod{async} & \cod{Deferred<T>} & resultado em paralelo \\
        \cod{runBlocking} & \cod{T} & ponte com código bloqueante \\
        \bottomrule
      \end{tabularx}
      \begin{erro}
        Usar \cod{runBlocking} no código do app trava a thread principal, que é justamente o que as coroutines evitam.
      \end{erro}
    \end{column}
  \end{columns}
\end{frame}

\begin{frame}[fragile]{Dispatchers e \cod{withContext}}
  \begin{columns}[T]
    \begin{column}{0.5\textwidth}
      \footnotesize
      \begin{tabularx}{\linewidth}{@{}l L@{}}
        \toprule
        \cabecalho{Dispatcher} & \cabecalho{Quando usar} \\
        \midrule
        \cod{Dispatchers.Main} & atualizar a UI; chamadas rápidas \\
        \cod{Dispatchers.IO} & rede, banco, arquivos (I/O bloqueante) \\
        \cod{Dispatchers.Default} & CPU: ordenar, filtrar, FFT, JSON grande \\
        \bottomrule
      \end{tabularx}
      \vspace{3pt}
      \begin{conceito}[Main-safety]
        Uma função é \emph{main-safe} quando pode ser chamada da thread principal com segurança, porque ela própria troca de dispatcher. Room, Retrofit e Firebase já são main-safe.
      \end{conceito}
    \end{column}
    \begin{column}{0.48\textwidth}
\begin{codigo}[title=Repositório main-safe]{Kotlin}
class LeituraRepository(
  private val io: CoroutineDispatcher = Dispatchers.IO
) {
  suspend fun exportarCsv(dados: List<Leitura>): File =
    withContext(io) {               // troca de thread
      val f = File(dir, "leituras.csv")
      f.writeText(dados.joinToString("\n") {
        "${it.instante};${it.valor}"
      })
      f
    }                               // volta ao chamador
}
\end{codigo}
      \begin{dica}
        Injetar o dispatcher permite trocá-lo nos testes (Módulo 14).
      \end{dica}
    \end{column}
  \end{columns}
\end{frame}

\begin{frame}{Concorrência estruturada}
  \begin{columns}[c]
    \begin{column}{0.5\textwidth}
      \centering
      \begin{tikzpicture}
        \node[caixaescura,minimum width=3.4cm] (s) at (0,0) {\cod{viewModelScope}};
        \node[caixa,minimum width=1.9cm] (j1) at (-1.6,-1.3) {\cod{launch} A};
        \node[caixa,minimum width=1.9cm] (j2) at (1.6,-1.3) {\cod{launch} B};
        \node[caixaambar,minimum width=1.5cm] (a1) at (0.7,-2.6) {\cod{async} B1};
        \node[caixaambar,minimum width=1.5cm] (a2) at (2.5,-2.6) {\cod{async} B2};
        \draw[seta] (s) -- (j1); \draw[seta] (s) -- (j2);
        \draw[seta] (j2) -- (a1); \draw[seta] (j2) -- (a2);
        \node[rotulo,text width=2.4cm] at (-1.6,-2.6) {cancelar o escopo cancela toda a árvore};
      \end{tikzpicture}
    \end{column}
    \begin{column}{0.48\textwidth}
      \small
      \begin{itemize}
        \item Toda coroutine nasce dentro de um \termo{CoroutineScope}.
        \item O pai espera os filhos terminarem.
        \item Cancelar o pai \textbf{cancela todos os filhos}.
        \item Uma exceção em um filho cancela o pai e os irmãos (exceto com \cod{SupervisorJob}).
      \end{itemize}
      \begin{conceito}[Resultado]
        Nenhuma coroutine fica ``solta'' consumindo bateria e rede depois que a tela que a criou é fechada.
      \end{conceito}
    \end{column}
  \end{columns}
\end{frame}

\begin{frame}[fragile]{Escopos no Android}
  \begin{columns}[T]
    \begin{column}{0.52\textwidth}
\begin{codigo}[title=Os três escopos do dia a dia]{Kotlin}
// 1) No ViewModel: vive enquanto o ViewModel viver
class PainelViewModel : ViewModel() {
  fun atualizar() = viewModelScope.launch {
    _estado.value = repo.carregar()
  }
}

// 2) Em Compose: ligado à composição
@Composable
fun BotaoSincronizar(onSync: suspend () -> Unit) {
  val scope = rememberCoroutineScope()
  Button(onClick = { scope.launch { onSync() } }) {
    Text("Sincronizar")
  }
}

// 3) Em Activity/Fragment: ligado ao ciclo de vida
lifecycleScope.launch { /* ... */ }
\end{codigo}
    \end{column}
    \begin{column}{0.46\textwidth}
      \footnotesize
      \begin{tabularx}{\linewidth}{@{}l L@{}}
        \toprule
        \cabecalho{Escopo} & \cabecalho{Cancelado quando} \\
        \midrule
        \cod{viewModelScope} & o ViewModel é destruído \\
        \cod{rememberCoroutineScope} & o composable sai da tela \\
        \cod{LaunchedEffect} & a chave muda ou sai da tela \\
        \cod{lifecycleScope} & a Activity é destruída \\
        \bottomrule
      \end{tabularx}
      \begin{erro}[Evite \cod{GlobalScope}]
        Coroutines em \cod{GlobalScope} não são canceladas automaticamente: causam vazamentos e trabalho desnecessário.
      \end{erro}
    \end{column}
  \end{columns}
\end{frame}

\begin{frame}[fragile]{Cancelamento e exceções}
  \begin{columns}[T]
    \begin{column}{0.5\textwidth}
\begin{codigo}[title=Cancelamento é cooperativo]{Kotlin}
val job = scope.launch(Dispatchers.Default) {
  for (amostra in amostras) {
    ensureActive()      // ou: if (!isActive) return
    processar(amostra)  // laço de CPU
  }
}
job.cancel()

// Timeout
val r = withTimeoutOrNull(3_000) {
  conectarSensor()
} ?: "sensor não respondeu"
\end{codigo}
      \begin{dica}
        Funções da biblioteca (\cod{delay}, \cod{withContext}, chamadas Room/Retrofit) já verificam o cancelamento.
      \end{dica}
    \end{column}
    \begin{column}{0.48\textwidth}
\begin{codigo}[title=Tratando falhas]{Kotlin}
viewModelScope.launch {
  try {
    val dados = repo.carregar()
    _estado.value = Estado.Sucesso(dados)
  } catch (e: CancellationException) {
    throw e           // nunca engula!
  } catch (e: IOException) {
    _estado.value = Estado.Erro("Sem conexão")
  }
}

// Filhos independentes: falha de um
// não cancela os outros
val escopo = CoroutineScope(
  SupervisorJob() + Dispatchers.Main
)
\end{codigo}
    \end{column}
  \end{columns}
\end{frame}

% =============================================================================
\section{Flow: fluxos de dados}

\begin{frame}[fragile]{O que é um \cod{Flow}}
  \begin{columns}[T]
    \begin{column}{0.5\textwidth}
      \small
      \begin{itemize}
        \item Uma \cod{suspend fun} devolve \textbf{um} valor. Um \termo{Flow} emite \textbf{vários valores ao longo do tempo}.
        \item É \termo{frio} (\emph{cold}): o código só roda quando alguém chama \cod{collect}, e roda de novo a cada coletor.
        \item Ideal para leituras de sensor, atualizações do banco e mensagens em tempo real.
      \end{itemize}
      \centering
      \begin{tikzpicture}[node distance=4mm]
        \node[caixaescura] (p) {produtor\\ \cod{emit()}};
        \node[caixaambar,right=of p] (o) {operadores\\ \cod{map}, \cod{filter}};
        \node[caixa,right=of o] (c) {coletor\\ \cod{collect}};
        \draw[seta] (p) -- (o); \draw[seta] (o) -- (c);
      \end{tikzpicture}
    \end{column}
    \begin{column}{0.48\textwidth}
\begin{codigo}[title=Simulando um sensor]{Kotlin}
fun temperaturaFlow(): Flow<Double> = flow {
  var t = 24.0
  while (true) {
    t += Random.nextDouble(-0.3, 0.3)
    emit(t)          // envia um valor
    delay(1_000)     // 1 Hz
  }
}

// consumindo
scope.launch {
  temperaturaFlow()
    .map { "%.1f °C".format(it) }
    .collect { texto -> println(texto) }
}
\end{codigo}
    \end{column}
  \end{columns}
\end{frame}

\begin{frame}[fragile]{Operadores úteis}
  \begin{columns}[T]
    \begin{column}{0.52\textwidth}
\begin{codigo}[title=Pipeline reativo]{Kotlin}
val alertas: Flow<String> = temperaturaFlow()
  .onEach { Log.d(TAG, "bruto: $it") }
  .map { it.arredondar(1) }
  .distinctUntilChanged()        // ignora repetidos
  .filter { it > 30.0 }
  .map { "Alerta: $it °C" }
  .catch { e -> emit("Erro: ${e.message}") }
  .flowOn(Dispatchers.Default)   // upstream em Default

// Combinar dois fluxos
val conforto = combine(tempFlow, umidFlow) { t, u ->
  indiceDeCalor(t, u)
}

// Busca com espera de digitação
consultaFlow
  .debounce(300)
  .flatMapLatest { termo -> repo.buscar(termo) }
\end{codigo}
    \end{column}
    \begin{column}{0.46\textwidth}
      \scriptsize
      \begin{tabularx}{\linewidth}{@{}l L@{}}
        \toprule
        \cabecalho{Operador} & \cabecalho{Efeito} \\
        \midrule
        \cod{map}, \cod{filter} & transformar e selecionar \\
        \cod{onEach} & efeito colateral (log) \\
        \cod{distinctUntilChanged} & descarta valores repetidos \\
        \cod{debounce(ms)} & espera o fluxo ``acalmar'' \\
        \cod{sample(ms)} & um valor a cada período \\
        \cod{combine} & junta os últimos valores de vários fluxos \\
        \cod{flatMapLatest} & troca de fluxo e cancela o anterior \\
        \cod{catch} & trata exceções do upstream \\
        \cod{flowOn} & muda o dispatcher do upstream \\
        \cod{take}, \cod{first} & limita a quantidade \\
        \bottomrule
      \end{tabularx}
    \end{column}
  \end{columns}
\end{frame}

\begin{frame}[fragile]{Convertendo callbacks em Flow: \cod{callbackFlow}}
  \begin{columns}[T]
    \begin{column}{0.56\textwidth}
\begin{codigo}[title=Sensor do aparelho como Flow]{Kotlin}
fun SensorManager.leituras(tipo: Int): Flow<FloatArray> =
  callbackFlow {
    val sensor = getDefaultSensor(tipo)
    val listener = object : SensorEventListener {
      override fun onSensorChanged(e: SensorEvent) {
        trySend(e.values.clone())   // envia ao Flow
      }
      override fun onAccuracyChanged(s: Sensor?, a: Int) {}
    }
    registerListener(listener, sensor,
      SensorManager.SENSOR_DELAY_UI)
    // obrigatório: limpa ao cancelar a coleta
    awaitClose { unregisterListener(listener) }
  }
\end{codigo}
    \end{column}
    \begin{column}{0.42\textwidth}
      \begin{conceito}[Quando usar]
        APIs baseadas em \emph{listeners}: sensores, localização, BLE, Firebase, MQTT. O \cod{callbackFlow} transforma tudo isso em \cod{Flow}, com cancelamento correto.
      \end{conceito}
      \begin{atencao}
        Esquecer o \cod{awaitClose} deixa o listener registrado, gastando bateria mesmo depois que ninguém mais observa. Voltaremos a esse padrão nos Módulos 12 e 13.
      \end{atencao}
    \end{column}
  \end{columns}
\end{frame}

% =============================================================================
\section{StateFlow e SharedFlow}

\begin{frame}[fragile]{\cod{StateFlow}: estado observável}
  \begin{columns}[T]
    \begin{column}{0.53\textwidth}
\begin{codigo}[title=Padrão do ViewModel]{Kotlin}
class PainelViewModel(
  private val repo: LeituraRepository
) : ViewModel() {
  // mutável e privado / somente leitura e público
  private val _ultima = MutableStateFlow<Double?>(null)
  val ultima: StateFlow<Double?> = _ultima.asStateFlow()

  fun simular() {
    _ultima.value = 25.3     // emite para os coletores
    _ultima.update { (it ?: 0.0) + 0.1 }  // atômico
  }

  // Converter um Flow frio em StateFlow
  val media: StateFlow<Double> = repo.leituras()
    .map { lista -> lista.map { it.valor }.average() }
    .stateIn(
      scope = viewModelScope,
      started = SharingStarted.WhileSubscribed(5_000),
      initialValue = 0.0
    )
}
\end{codigo}
    \end{column}
    \begin{column}{0.45\textwidth}
      \begin{conceito}[StateFlow]
        Fluxo \textbf{quente} que sempre tem um valor atual (\cod{.value}). Novos coletores recebem imediatamente o último valor. Valores iguais em sequência não são reemitidos.
      \end{conceito}
      \begin{dica}[\cod{WhileSubscribed(5\_000)}]
        Mantém o fluxo ativo por 5\,s após o último coletor sair. Assim uma rotação de tela não reinicia a consulta, mas sair da tela para tudo.
      \end{dica}
      \begin{conceito}[Na UI (Módulo 07)]
        \cod{val v by vm.ultima}\\\cod{\ \ .collectAsStateWithLifecycle()}
      \end{conceito}
    \end{column}
  \end{columns}
\end{frame}

\begin{frame}[fragile]{\cod{SharedFlow} e comparação}
  \begin{columns}[T]
    \begin{column}{0.48\textwidth}
\begin{codigo}[title=Eventos de uso único]{Kotlin}
private val _eventos = MutableSharedFlow<String>()
val eventos = _eventos.asSharedFlow()

fun salvar() = viewModelScope.launch {
  repo.salvar()
  _eventos.emit("Leitura salva!")   // ex.: Snackbar
}
\end{codigo}
      \begin{dica}
        A recomendação atual é modelar até mensagens como parte do estado (ex.: \cod{mensagem: String?} no \emph{UiState}) e limpá-las depois de exibidas. Use \cod{SharedFlow} quando o evento realmente não for estado.
      \end{dica}
    \end{column}
    \begin{column}{0.50\textwidth}
      \scriptsize
      \begin{tabularx}{\linewidth}{@{}l L L L@{}}
        \toprule
        & \cabecalho{Flow} & \cabecalho{StateFlow} & \cabecalho{SharedFlow} \\
        \midrule
        Tipo & frio & quente & quente \\
        Valor atual & não & sim (\cod{.value}) & não (ou \emph{replay}) \\
        Valor inicial & --- & obrigatório & --- \\
        Coletores & cada um roda o produtor & compartilham & compartilham \\
        Uso típico & dados do repositório & estado da tela & eventos \\
        \bottomrule
      \end{tabularx}
    \end{column}
  \end{columns}
\end{frame}

% =============================================================================
\section{Fechamento}

\begin{frame}{Resumo}
  \begin{resumo}
    \begin{itemize}
      \item A thread principal só desenha a UI; I/O e CPU pesada vão para \cod{Dispatchers.IO}/\cod{Default}.
      \item \cod{suspend} marca funções que podem suspender; \cod{launch} e \cod{async} criam coroutines.
      \item Use os escopos com ciclo de vida (\cod{viewModelScope}, \cod{rememberCoroutineScope}) e nunca \cod{GlobalScope}.
      \item Cancelamento é cooperativo; relance \cod{CancellationException}.
      \item \cod{Flow} representa vários valores no tempo; operadores montam pipelines reativos.
      \item \cod{StateFlow} expõe o estado da tela; \cod{stateIn} converte fluxos frios.
    \end{itemize}
  \end{resumo}
\end{frame}

\begin{frame}{Glossário (1/2)}
  \begin{glossario}
    \emph{Thread} & linha de execução dentro de um processo, com pilha própria \\
    \emph{Thread} principal & desenha a interface e trata os toques do usuário \\
    ANR & diálogo exibido quando a \emph{thread} principal deixa de responder \\
    Concorrência & tarefas em andamento no mesmo período, que se alternam \\
    Paralelismo & tarefas executando no mesmo instante, em núcleos diferentes \\
    \emph{Coroutine} & computação que pode ser suspensa e retomada sem bloquear a \emph{thread} \\
    \cod{suspend} & marca uma função que pode suspender a coroutine \\
    \emph{Dispatcher} & define em que \emph{threads} a coroutine roda (\cod{Main}, \cod{IO}, \cod{Default}) \\
    Concorrência estruturada & tarefas presas a um escopo: cancelar o escopo cancela as filhas \\
  \end{glossario}
\end{frame}

\begin{frame}{Glossário (2/2)}
  \begin{glossario}
    Cancelamento cooperativo & a coroutine para em pontos de suspensão ou ao verificar \cod{isActive} \\
    \cod{Flow} & fluxo frio de vários valores ao longo do tempo \\
    \cod{StateFlow} & fluxo quente que sempre guarda o valor atual \\
    \cod{SharedFlow} & fluxo quente para eventos, com vários coletores \\
    Condição de corrida & resultado que depende da ordem imprevisível de acesso a um dado \\
    \cod{Mutex} & trava que garante exclusão mútua entre coroutines \\
    Contrapressão & estratégia para quando o produtor é mais rápido que o consumidor \\
  \end{glossario}
\end{frame}

\begin{frame}{Questões de revisão}
  \begin{questoes}
    \item Por que uma requisição de rede não pode rodar na \emph{thread} principal? O que o usuário veria?
    \item Diferencie concorrência de paralelismo com um exemplo no celular.
    \item Uma tarefa de I/O e uma de CPU merecem o mesmo tratamento? Que \emph{dispatcher} usar em cada caso?
    \item O que acontece com a \emph{thread} quando uma função \cod{suspend} chega a um ponto de suspensão?
    \item Explique a concorrência estruturada e o que ela evita quando a tela é fechada.
    \item Duas coroutines fazem \cod{contador++} mil vezes cada no \cod{Dispatchers.Default}. Por que o total pode ficar abaixo de 2000? Como corrigir?
    \item Qual a diferença entre um \cod{Flow} frio e um \cod{StateFlow}?
    \item Um sensor emite 100 leituras por segundo e a tela só precisa da última. Que estratégia de contrapressão usar?
  \end{questoes}
\end{frame}

\begin{frame}{Exercícios}
  \begin{exercicio}[Exercício 1 --- Paralelismo]
    Escreva três funções \cod{suspend} que simulam leituras com \cod{delay} de 1\,s, 2\,s e 3\,s. Meça o tempo total chamando-as em sequência e depois com \cod{async}. Explique a diferença.
  \end{exercicio}
  \begin{exercicio}[Exercício 2 --- Flow]
    Crie um \cod{Flow<Int>} que simula um ADC de 12 bits a 10\,Hz. Monte um pipeline que converte para volts, calcula a média de janelas de 10 amostras e emite alerta acima de 3,0\,V.
  \end{exercicio}
  \begin{exercicio}[Exercício 3 --- Cancelamento]
    Inicie um laço de CPU com \cod{launch(Dispatchers.Default)} sem \cod{ensureActive()} e cancele após 100\,ms. O que acontece? Corrija.
  \end{exercicio}
  \begin{projeto}[LabMonitor --- Etapa 6]
    Implemente \cod{SimuladorSensor} com um \cod{Flow<Leitura>} por bancada e um \cod{PainelViewModel} que expõe a última leitura e a média dos últimos 60\,s via \cod{StateFlow}.
  \end{projeto}
\end{frame}

\begin{frame}{Para aprofundar}
  \small
  \begin{itemize}
    \item Coroutines no Android: \url{https://developer.android.com/kotlin/coroutines}
    \item Boas práticas com coroutines: \url{https://developer.android.com/kotlin/coroutines/coroutines-best-practices}
    \item Guia oficial do Kotlin: \url{https://kotlinlang.org/docs/coroutines-guide.html}
    \item Flow no Android: \url{https://developer.android.com/kotlin/flow}
    \item StateFlow e SharedFlow: \url{https://developer.android.com/kotlin/flow/stateflow-and-sharedflow}
  \end{itemize}
\end{frame}

\begin{frame}{Respostas das questões de revisão}
  \begin{respostas}
    \item A \emph{thread} principal desenha a tela e trata toques; bloqueá-la congela a interface, perde quadros e, depois de cerca de 5\,s, gera um ANR. O Android ainda lança \cod{NetworkOnMainThreadException}.
    \item Concorrência: um download e uma animação se alternando em um núcleo. Paralelismo: duas imagens decodificadas ao mesmo tempo em dois núcleos.
    \item Não. I/O passa quase todo o tempo esperando (\cod{Dispatchers.IO} ou APIs que já suspendem); CPU ocupa o núcleo (\cod{Dispatchers.Default}).
    \item A coroutine pausa e a \emph{thread} fica livre para outras tarefas; quando o resultado chega, ela continua de onde parou, talvez em outra \emph{thread} do mesmo \emph{dispatcher}.
    \item Toda coroutine pertence a um escopo e não vive mais que ele. Ao fechar a tela, o escopo (\cod{viewModelScope}, \cod{rememberCoroutineScope}) cancela as filhas: nada de bateria, rede e memória desperdiçadas.
    \item \cod{contador++} é ler, somar e gravar: em \emph{threads} diferentes, as duas podem ler o mesmo valor e uma atualização se perde. Use \cod{AtomicInteger}, \cod{Mutex} ou confine o contador a uma \emph{thread}.
    \item O \cod{Flow} frio só produz quando coletado e recomeça para cada coletor; o \cod{StateFlow} é quente, sempre tem um valor e entrega o último a quem chega.
    \item Manter só o mais recente: \cod{conflate()} ou \cod{collectLatest}. Um \cod{StateFlow} já se comporta assim.
  \end{respostas}
\end{frame}

\encerramento{Módulo 07 --- Estado, navegação e animações: remember, state hoisting, efeitos, ciclo de vida, Navigation 3 e animações.}

\end{document}
